\chapter{Trabalhos Relacionados}
\label{cap:trabalhos}

A literatura sobre transparência pública e a literatura sobre Data Warehouse raramente conversam entre si. A primeira concentra-se em avaliação de portais, marcos legais e teoria do controle social. A segunda, em arquitetura de dados, modelagem dimensional e ferramentas de \textit{Business Intelligence}. Este capítulo percorre as duas tradições e identifica o espaço que esta pesquisa pretende ocupar.

\section{Avaliação de Portais de Transparência Municipais}

O esforço de medir a qualidade dos portais de transparência brasileiros não é novo. \citeonline{melo2019proposta} propõem um índice bidimensional de transparência da informação público-eletrônica, dividido entre transparência ativa (o que o portal disponibiliza espontaneamente) e transparência passiva (a resposta a pedidos formais de informação). Os autores aplicam o índice a municípios brasileiros e mostram que a maior parte deles atinge níveis baixos a moderados, com forte heterogeneidade entre regiões. O resultado importa porque revela uma assimetria estrutural: o cidadão de um município pequeno do interior tem, em média, menos acesso à informação do que o cidadão de uma capital. A LAI é a mesma; a implementação não.

\citeonline{marco2022transparencia} adotam uma abordagem complementar e investigam a transparência municipal pela perspectiva dos Observatórios Sociais. Em vez de medir o portal pelo lado da oferta, os autores ouvem o lado da demanda, entrevistando organizações civis que efetivamente usam essas plataformas para fiscalização. O achado é consistente com o argumento freireano: mesmo grupos organizados, com conhecimento técnico médio, encontram dificuldades para cruzar dados de despesas, licitações e contratos. A barreira de uso permanece, ainda quando o usuário é qualificado. Se o cidadão organizado tropeça, o cidadão comum sequer começa a caminhar.

\citeonline{pinho2009accountability} oferecem o pano de fundo teórico para essas avaliações empíricas. Ao discutir se ``\textit{accountability}'' tem tradução adequada para o português, os autores expõem a distância entre publicar informação e prestar contas de fato. Transparência é condição necessária, mas não suficiente, para \textit{accountability} democrática. Esse referencial dá nome ao problema que esta pesquisa endereça e justifica por que a simples publicação de dados em portais não fecha o ciclo da participação cidadã.

Em registro complementar, \citeonline{souza2025transparencia} reúnem evidência recente sobre transparência e controle social no contexto do orçamento participativo brasileiro. Os autores identificam três obstáculos persistentes à participação qualificada: baixa escolaridade política, descontinuidade administrativa entre gestões e apropriação estratégica da participação pelo próprio poder público, fenômeno em que mecanismos formalmente participativos são usados para legitimar decisões já tomadas. A leitura conjunta com os trabalhos anteriores mostra que o problema da transparência efetiva não tem se resolvido com o passar do tempo, apesar dos avanços legais e do amadurecimento dos portais.

\section{Data Warehouse Aplicado a Dados Governamentais}

A aplicação de Data Warehouse a dados públicos é menos documentada na literatura acadêmica brasileira do que se poderia esperar. A obra de referência de \citeonline{kimball2013data} oferece a metodologia, mas trata de casos majoritariamente empresariais. As propriedades que definem o DW (orientação a assunto, integração, não volatilidade e variação no tempo) traduzem-se sem dificuldade para o contexto orçamentário, porém os exemplos clássicos da bibliografia raramente partem de dados públicos.

Na prática, projetos internacionais como o Open Spending, mantido pela Open Knowledge Foundation, mostraram nas duas últimas décadas a viabilidade de consolidar orçamentos públicos em estruturas analíticas reutilizáveis. O projeto agrega dados orçamentários de centenas de instituições em formato padronizado, oferecendo visualizações e API. A escolha de modelo de dados, no entanto, é genérica e não captura especificidades da contabilidade pública brasileira (PPA, LDO, LOA, classificação MCASP). Para usuários brasileiros, isso significa que o Open Spending é referência arquitetural útil, mas exige adaptação substancial para responder às perguntas que importam aqui.

A literatura internacional sobre o tema introduz, ainda, o conceito de \textit{citizens' budget} como prática complementar ao DW. \citeonline{bilge2015citizens} descreve esse instrumento como uma versão simplificada e explicativa do orçamento público, voltada ao cidadão não especialista, e cataloga sua adoção em países como Estados Unidos, Reino Unido, África do Sul e Filipinas. O ponto importante para esta pesquisa não é o documento simplificado em si, mas o reconhecimento, em escala internacional, de que a intelegibilidade orçamentária é um problema de \textit{design} de informação, e não apenas de publicação. A solução proposta aqui via DW se distingue por permitir interação analítica em vez de leitura passiva de um relatório resumido.

No Brasil, o Painel do Orçamento Federal, mantido pelo Tesouro Nacional, e o Portal da Transparência da Controladoria-Geral da União representam o esforço público mais maduro nessa direção. Ambos consolidam dados em estruturas analíticas e oferecem visualizações temáticas. Esses sistemas, entretanto, cobrem apenas a esfera federal. A execução orçamentária municipal segue espalhada por portais individuais, com qualidade e profundidade muito variáveis. Não há equivalente municipal consolidado e a iniciativa de construir um para um município específico permanece responsabilidade local.

\comment[inline]{Pesquisar e adicionar referências acadêmicas específicas sobre DW aplicado a dados governamentais no Brasil. Termos de busca sugeridos: ``data warehouse setor público brasil'', ``dimensional model government data'', ``transparência fiscal modelo dimensional''. Verificar trabalhos de pós-graduação em Administração Pública e em Ciência da Computação que cruzem os temas.}

\section{Visualização e \textit{Civic Tech}}

Iniciativas de \textit{civic tech} brasileiras complementam o quadro institucional. A Operação Serenata de Amor, por exemplo, partiu de uma abordagem oposta à dos órgãos oficiais. Em vez de tentar cobrir tudo, focou em um problema específico (verificação de despesas da Cota para Exercício da Atividade Parlamentar) e construiu sobre ele algoritmos e visualizações que tornaram públicos casos concretos de irregularidade. A lição metodológica é relevante para este trabalho: foco temático bem definido produz mais impacto cívico do que generalidade técnica.

Esse tipo de iniciativa costuma ter vida útil ligada à dedicação de voluntários e a financiamentos pontuais. Sistemas mantidos institucionalmente são mais duráveis, mas costumam ser conservadores na escolha do recorte. Há aí uma tensão produtiva entre o experimento de \textit{civic tech} e o sistema institucional, e o presente trabalho se posiciona em um terreno intermediário: rigor metodológico de DW combinado com recorte temático claro de uma cidade específica.

\section{Posicionamento desta Pesquisa}

Em relação aos trabalhos discutidos, esta pesquisa se diferencia em quatro pontos. Primeiro, o recorte municipal específico: Juiz de Fora, com a granularidade dos dados realmente publicados pela prefeitura. Segundo, a integração entre receitas e despesas em um único modelo dimensional, raramente feita em sistemas existentes. Terceiro, a documentação explícita das inconsistências encontradas nas fontes, tratada como achado de pesquisa e não como ruído a ser silenciosamente corrigido. Quarto, a articulação teórica com a tradição freireana de educação libertadora, que ancora a justificativa do projeto em algo além da eficiência técnica.

A Tabela~\ref{tab:posicionamento} resume essa comparação.

\begin{table}[htb]
\caption{Posicionamento desta pesquisa frente aos trabalhos relacionados}
\label{tab:posicionamento}
\centering
\small
\begin{tabular}{p{3.2cm} p{2.0cm} p{2.0cm} p{2.0cm} p{2.0cm}}
\hline
\textbf{Característica} & \textbf{Índices de transparência} & \textbf{Painel CGU/Tesouro} & \textbf{\textit{Civic tech}} & \textbf{Este trabalho} \\
\hline
Foco municipal & Sim & Não & Variável & Sim \\
Modelo dimensional explícito & Não & Parcial & Não & Sim \\
Integração receita+despesa & Não & Sim & Não & Sim \\
Documentação de inconsistências & Não & Não & Parcial & Sim \\
Fundamentação freireana & Não & Não & Não & Sim \\
\hline
\end{tabular}
\end{table}

O capítulo seguinte detalha a metodologia adotada para concretizar essas diferenças no caso de Juiz de Fora.
